% Session 7 -- DMRG: differentiating the network erases one tensor and leaves a
% small eigenvalue problem; sweeping that problem along the chain is the algorithm.
% Placement: sec. "Finding the ground state: the variational principle and DMRG",
% after the "erasing that tensor from the diagram" paragraph.
%
% LAYOUT NOTE: pure TikZ, canvas coordinates; one-off tensor styles are \tikzset
% inside this picture only. Panel titles are short tags, with the explanatory line
% placed under each panel, so nothing collides across the top of the figure.
\begin{aqmfigure}{fig:s07:dmrg}{How DMRG turns an exponentially large minimization into a sequence of
small ones. (a) The numerator of the Rayleigh quotient (7.11) is a single closed network, made of the
ket MPS, the Hamiltonian as an MPO and the bra MPS, all contracted together. It has no free legs, so
it is a number. (b) Every tensor enters that network once in the bra layer and once in the ket layer, so
the energy is a quadratic form in each tensor and linear in its bra copy, and the derivative with
respect to the bra tensor at site $k$ removes it. What is left has
exactly three open legs, the two virtual bonds and the physical index of site $k$, which are the
indices of the tensor $M_k^{s_k}$ that we want to find. Everything to the left and right of $k$ has
been reduced to an environment, and the stationarity condition is a generalized eigenvalue problem,
which is an ordinary one in an orthonormal gauge, with the size of a single tensor. (c) We solve it, move one site along, solve again, and sweep back and
forth until the energy stops decreasing. This is the density matrix renormalization group, and at no
point do we build an object of size $2^L$.}

\begin{tikzpicture}[x=1cm, y=1cm]
\tikzset{
  tens/.style={aqm, circle, draw=black, fill=aqmfill, minimum size=6.5pt, inner sep=0pt},
  tensop/.style={aqm, rectangle, draw=black, fill=white, minimum size=6.5pt, inner sep=0pt},
  leg/.style={aqm, draw=black, line width=\aqmthinwidth},
  openleg/.style={aqm, draw=aqmbroken, line width=1.15pt}
}

% ============================ (a) the closed network ========================
\node[aqmpanel, anchor=west] at (-0.05,0.60) {(a)};

\node[aqmlabsmall, anchor=east] at (0.42,0.10)  {$\langle\Psi|$};
\node[aqmlabsmall, anchor=east] at (0.42,-0.52) {$H$};
\node[aqmlabsmall, anchor=east] at (0.42,-1.14) {$|\Psi\rangle$};

\foreach \i in {0,1,2,3,4} {
  \node[tens]   (b\i) at (0.62+0.58*\i, 0.10)  {};
  \node[tensop] (h\i) at (0.62+0.58*\i,-0.52)  {};
  \node[tens]   (k\i) at (0.62+0.58*\i,-1.14)  {};
  \draw[leg] (b\i) -- (h\i);
  \draw[leg] (h\i) -- (k\i);
}
\foreach \i/\j in {0/1, 1/2, 2/3, 3/4} {
  \draw[leg] (b\i) -- (b\j);
  \draw[leg] (h\i) -- (h\j);
  \draw[leg] (k\i) -- (k\j);
}
\node[aqmlabsmall, anchor=north, text=aqmmuted, align=center] at (1.78,-1.44)
  {$\langle\Psi|H|\Psi\rangle$: no free legs,\\[-1pt] so the contraction is a number};

% ============================ the differentiation arrow =====================
\draw[aqmarrowaccent] (3.40,-0.52) -- (4.92,-0.52);
\node[aqmlabsmall, anchor=south, align=center, text=aqmaccent] at (4.15,-0.44)
  {erase the bra\\[-1pt] tensor at $k$};

% ============================ (b) the hole ==================================
\node[aqmpanel, anchor=west] at (5.25,0.60) {(b)};

\foreach \i in {0,1,2,3,4} {
  \node[tensop] (H\i) at (5.65+0.58*\i,-0.52)  {};
  \node[tens]   (K\i) at (5.65+0.58*\i,-1.14)  {};
  \draw[leg] (H\i) -- (K\i);
}
\foreach \i in {0,1,3,4} {
  \node[tens] (B\i) at (5.65+0.58*\i, 0.10)  {};
  \draw[leg] (B\i) -- (H\i);
}
\foreach \i/\j in {0/1, 1/2, 2/3, 3/4} {
  \draw[leg] (H\i) -- (H\j);
  \draw[leg] (K\i) -- (K\j);
}
% the hole at site k, and the three legs it leaves open
\draw[openleg] (B1) -- (6.60,0.10);
\draw[openleg] (7.00,0.10) -- (B3);
\draw[openleg] (H2) -- (6.80,-0.12);
\node[circle, draw=aqmbroken, dashed, line width=\aqmthinwidth, fill=none,
      minimum size=12pt, inner sep=0pt] at (6.80,0.10) {};
\node[aqmlabsmall, anchor=south, align=center, text=aqmbroken] at (6.80,0.34)
  {site $k$ erased};
\node[aqmlabsmall, anchor=north, text=aqmbroken, align=center] at (6.80,-1.44)
  {three open legs: exactly the\\[-1pt] indices of one tensor $M_k^{s_k}$};

\node[aqmboxaccent, anchor=west, text width=4.5cm, align=center] at (8.75,-0.52)
  {$H^{\text{eff}}_k\,M_k^{s_k} = \lambda\,M_k^{s_k}$\\[3pt]
   {\footnotesize\textcolor{aqmmuted}{one small eigenvalue problem; the rest of the chain is
   now just an environment}}};

% ============================ (c) the sweep =================================
\node[aqmpanel, anchor=west] at (-0.05,-2.48) {(c)};

\foreach \i in {0,...,7} {
  \node[tens] (s\i) at (0.62+0.58*\i,-3.20) {};
  \draw[leg] (0.62+0.58*\i,-3.28) -- (0.62+0.58*\i,-3.52);
}
\foreach \i/\j in {0/1, 1/2, 2/3, 3/4, 4/5, 5/6, 6/7} { \draw[leg] (s\i) -- (s\j); }

\draw[aqmarrowaccent] (0.50,-2.86) -- (4.82,-2.86);
\node[aqmlabsmall, anchor=west, text=aqmaccent] at (5.05,-2.86)
  {solve one tensor, move one site along};
\draw[aqmarrowmuted] (4.82,-3.78) -- (0.50,-3.78);
\node[aqmlabsmall, anchor=west, text=aqmmuted] at (5.05,-3.78)
  {\dots{} then sweep back, until the energy stops decreasing};

\end{tikzpicture}

\end{aqmfigure}
